\section{在对数域中稳定计算概率}
\label{sec:08-Numerical-Analysis/07-Numerical-Methods-for-ML/01}

一个词性标注模型跑通了，在几十个词的句子上一切正常。换成一整篇新闻稿——一千多个词——序列似然返回 0，困惑度打印出 \texttt{inf}，评估脚本停在那里。

代码没写错。序列似然是每一步条件概率的连乘，每一项大约 \(10^{-8}\)。乘到第三十九项，乘积落到 \(10^{-312}\) 附近，已经掉进 float64 的次正规区（约 \(2.2\times10^{-308}\) 以下），有效位开始一位一位地少；第四十一项之后连次正规都装不下，乘积变成精确的零。后面九百多个因子乘的都是零，一个也没参与。

换更宽的类型也不解决问题。float128 把下界从 \(10^{-308}\) 推到 \(10^{-4932}\)，序列长到六百词又是零。真实答案 \(10^{-8000}\) 不是一个``很小的数''，它是一个 float64 从设计上就不打算表示的数。要绕开它，不能靠加位数，只能换一个表示这个量的方式。

\subsection{连乘变连加，范围问题就此消失}

对数正是那个变换。取 \(\ell=\sum_t\log p_t\)，一千项每项约 \(-18.4\)，加起来是 \(-18420\)——一个再普通不过的双精度数，离溢出边界还有八个数量级。而且它不会随序列变长而恶化：加法只让结果线性地变大，不像乘法那样把指数一路往下堆。

对数域还顺带修好了另一件事。概率的相对精度在乘法里会累积，\(n\) 次乘法的相对误差大致是 \(nu\)；换成加法后，误差是绝对量级的 \(nu\abs{\ell}\)，对数似然本身有多大，误差就按比例来——这恰恰是我们要用它做的事（比较两个模型的对数似然差）所需要的精度。

到这里似乎收工了。但概率模型很少只做连乘。

\subsection{对数域里的求和是新的麻烦}

隐 Markov 模型的前向递推每走一步都要对上一时刻的隐状态求和，\(\alpha_t(j)=\sum_i\alpha_{t-1}(i)a_{ij}b_j(o_t)\)；混合模型的 E 步要算各分量的加权和；Bayes 更新要算归一化常数。乘法在对数域里有对应物，加法没有。已知的是 \(x_i=\log u_i\)，要的是

\[
\log\sum_i u_i=\log\sum_i e^{x_i}.
\]

照字面写成程序就是三步：取指数、求和、取对数。而 \(x_i\) 恰恰是刚才那批很负的数。\(e^{-18420}\) 在 float64 里是零，一堆零加起来还是零，取对数得到 \(-\infty\)——绕了一圈回到原地。另一头也堵着：分类模型最后一层的 logit 若是 \(1000\)，\(e^{1000}\approx2\times10^{434}\) 溢出为 \(+\infty\)，随后 \(\infty/\infty\) 给出 NaN。指数函数把浮点数那个跨越六百多个数量级的宽阔区间，压缩成输入侧一条只有一千四百多宽的窄缝。

\subsection{提出最大值，恒等式没动，窗口跟着输入走}

出路是一行代数。令 \(m=\max_i x_i\)，把 \(e^m\) 从和式里提出来：

\[
\logsumexp(x)=\log\sum_i e^{x_i}
=\log\Big(e^{m}\sum_i e^{x_i-m}\Big)
=m+\log\sum_i e^{x_i-m}.
\]

在实数上这是恒等式，无条件成立。浮点上它还要两个条件：\(m\) 得是有限的，以及和式里至少有一项不下溢。后一条自动满足——最大项减去自己恰好是 \(e^0=1\)，这是世界上最安全的指数计算。前一条不自动满足，本节末尾会回来。

平移之后每一项都落在 \((0,1]\)，和落在 \([1,n]\)，取对数落在 \([0,\log n]\)，最后加回 \(m\)。整条路径上没有一个中间量离开浮点数的舒适区，而 \(m\) 是多大都无所谓——它从头到尾只参加加减法。

\begin{center}
\begin{tikzpicture}[x=1cm,y=1cm,font=\small,>=stealth]
  \fill[red!16] (1,2.5) rectangle (3.27,2.95);
  \fill[green!22] (3.27,2.5) rectangle (10.55,2.95);
  \fill[red!16] (10.55,2.5) rectangle (13,2.95);
  \draw[black!45] (1,2.5) rectangle (13,2.95);
  \node[above,font=\footnotesize] at (7,3.02) {直接取指数再求和};
  \node[font=\footnotesize,black!70] at (2.1,2.72) {全下溢};
  \node[font=\footnotesize,black!70] at (11.75,2.72) {上溢};
  \node[font=\footnotesize] at (6.9,2.72) {可用};
  \fill[green!22] (1,1.35) rectangle (13,1.8);
  \draw[black!45] (1,1.35) rectangle (13,1.8);
  \node[above,font=\footnotesize] at (7,1.87) {先减去最大值（LSE）};
  \node[font=\footnotesize] at (7,1.57) {全程可用};
  \draw[->,black!55] (0.7,0.9) -- (13.5,0.9);
  \foreach \xa/\lab in {2/{\(-1000\)},3.27/{\(-745\)},7/{\(0\)},10.55/{\(709\)},12/{\(1000\)}}{
    \draw[black!45] (\xa,0.8) -- (\xa,1.0);
    \node[below,font=\footnotesize,black!70] at (\xa,0.8) {\lab};
  }
  \node[font=\footnotesize,black!70] at (7,0.05) {输入分量 \(x_i\) 的量级（float64）};
\end{tikzpicture}

\emph{直接计算只在中间那段宽约 1450 的窗口里有效，越出左端全部下溢、结果是 \(-\infty\)，越出右端上溢、结果是 NaN。减去最大值之后，窗口不再钉在原点，而是跟着 \(\max_i x_i\) 平移——只要它本身是个合法浮点数，结果就合法。}
\end{center}

图里两条带子的宽度其实一样：能被安全指数化的动态范围是浮点格式给死的，谁也变不出更多。变的是这段窗口摆在哪儿。直译的写法把它焊在原点上，于是输入必须自己走进窗口；LSE 把窗口搬到输入身边。同一个数学表达式有无穷多种等价写法，它们的数值行为可以完全不同——这句话是本章三篇共用的判断，log-sum-exp 是它最干净的一个例子。

顺着这个分解还能读出一个不等式。\(1\le\sum_i e^{x_i-m}\le n\)，三边取对数再加回 \(m\)，得到

\[
\max_i x_i\le\logsumexp(x)\le\max_i x_i+\log n.
\]

LSE 因此是最大值的一个光滑版本：某一项远远领先时它几乎就等于那一项，几项接近时那不超过 \(\log n\) 的余量记下了竞争者的贡献。它的梯度更值得记住，

\[
\frac{\partial\logsumexp(x)}{\partial x_i}=\frac{e^{x_i}}{\sum_j e^{x_j}}=\softmax(x)_i,
\]

softmax 就是 LSE 对输入的敏感度。凸优化里的``函数值与次梯度''、统计里的``配分函数与充分统计量期望''是同一种配对结构，见\hyperref[sec:13-Machine-Learning/06-Probabilistic-Models/04]{``EM 用下界交替完成隐变量学习''}中对边缘似然的处理。

\subsection{要对数就别绕道概率}

分类训练要的是正确类别的负对数概率 \(-\log p_y\)。顺手的写法是先 softmax、后取 log，两步分开。这一绕，中间就多出了一整张概率表，而表里的极小项可能已经下溢为零；对零取对数是 \(-\infty\)，接着整个 batch 的梯度是 NaN。

直接由 logit 算就没有这一步：

\[
-\log p_y=\logsumexp(x)-x_y.
\]

即使 \(p_y\) 小到浮点数认不出它和零的区别，\(x_y-\logsumexp(x)\) 仍然是个规规矩矩的负数，比如 \(-8000\)——训练照样能从这个数上拿到有意义的梯度。同理，log-softmax 应当写成 \(\log p_i=x_i-\logsumexp(x)\)，而不是先算 \(\softmax\) 再取对数。PyTorch 的 \texttt{log\_softmax}、\texttt{cross\_entropy} 和 \texttt{logsumexp} 把这些步骤融在一个核里，理由既是省一趟内存，也是省掉那张会下溢的中间表。

同样的原则一路用到底：HMM 前向递推整条留在对数域，每一步的求和用 LSE；混合模型的 E 步算责任度时，先在对数域上算 \(\log\pi_k+\log\mathcal{N}(x\given\mu_k,\Sigma_k)\)，再对 \(k\) 做一次 LSE 归一化，中途不生成任何概率。详见\hyperref[sec:13-Machine-Learning/06-Probabilistic-Models/05]{``HMM 把隐变量排列成时间链''}与\hyperref[sec:06-Random-Processes/03-Markov-Chains/01]{``Markov 性质、转移矩阵与多步演化''}。

\subsection{下溢不总是灾难，不知道自己丢了什么才是}

减完最大值以后，仍然会有项下溢为零。输入 \((0,-1000)\) 在 float32 下算出概率 \((1,0)\)，第二个概率的真值约 \(e^{-1000}\)，它对归一化和的贡献本来就在舍入误差以下，丢了不影响任何后续计算。下溢本身不是 bug，在不该下溢的地方下溢才是。

真正会咬人的是全 mask 那一行。序列模型把填充位置的 logit 设成 \(-\infty\)，让 softmax 自动把它们压成零，这没问题；可如果某一行的全部位置都被遮住，\(m=-\infty\)，平移那一步出现 \(-\infty-(-\infty)\)，按 IEEE 754 就是 NaN。这正是前面挂账的那个条件。它不是数值技巧能修的——一整行都无效时不存在唯一正确的概率分布，跳过、报错、还是填成均匀分布，取决于这一行在你的任务里意味着什么。稳定实现能做的是把选项摆出来，不能替调用者做决定。

混合精度里还有一处容易漏。输入 logit 留在 fp16 没关系，但求最大值、指数求和与损失累加应当在 fp32 里做。减最大值解决的是动态范围，它不会把 fp16 已经丢掉的有效位还回来，这一点在\hyperref[sec:08-Numerical-Analysis/01-Floating-Point-and-Error/02]{``误差如何产生与放大''}讨论长求和时已经出现过。

\subsection{用数学性质当测试，比抄期望输出管用}

验证一个稳定实现，最省事也最没用的办法是抄几组参考输出——抄错了也看不出来，而且它只覆盖被抄到的那几个点。更有力的做法是拿定义里的恒等式去测：所有 logit 加同一个常数，softmax 的输出应当一位不变，LSE 应当恰好增加那个常数；每行概率非负且和为 \(1\)（容许舍入量级的偏差）；\(10^4\)、\(-10^4\)、全等输入、只剩一个有效位置的 mask，都应返回有限值而不是 Inf 或 NaN；log-softmax 的结果应与 \(x_i-\logsumexp(x)\) 逐元素一致。

平移不变性这一条尤其能干活。轴选错、广播方向反了、mask 加在了减最大值之后——这些错误在随手挑的例子上常常碰巧给出看着合理的数，却几乎不可能同时满足平移不变。一条来自定义的性质，抵得过一百个硬编码的用例。

改写表达式能救回范围，救不回信任：一段自己写的求导代码是对是错，恒等式检验不了。下一节把这件事交给两条独立的计算路径。
